\section{Compute Lifecycle: preentrenamiento, post-training e inference}

La clase rechaza la narrativa binaria de que «el preentrenamiento ya terminó». El preentrenamiento absorbe con eficiencia datos amplios y establece la base capability; el post-training enfoca la conducta, las herramientas y los objetivos de seguridad; el inference compute resuelve tareas concretas mediante context, search, reasoning o muestreo múltiple. Los tres difieren en rendimiento marginal, latency, cost y controlabilidad.

\subsection{Cada etapa resuelve un problema distinto}

El preentrenamiento conviene para capacidades amplias cuyo costo se amortiza entre todas las solicitudes; el post-training conviene para la conducta de producto o de políticas; la inference conviene para invertir de forma dinámica según la tarea. Si al base model le faltan conocimientos o representation, aumentar los inference tokens puede equivaler solo a adivinar durante más tiempo; si el post-training es insuficiente, incluso un base model potente puede no respetar el contract de las herramientas; si el serving budget es demasiado bajo, el reasoning complejo tampoco puede llevarse a la práctica.

\lecturefigure{15-compute-lifecycle.png}{El preentrenamiento, el post-training, el inference compute y el product feedback forman un ciclo de vida complementario.}{Subtítulos locales 34:40--36:50; rediseño conceptual.}

\begin{knowledgebox}{Lectura de la figura: comparar el total cost, no solo los FLOPs de entrenamiento}
La economía de un modelo incluye la amortización del entrenamiento, la post-training iteration, el accelerator/time por solicitud, el context storage, la cache, las tool calls y los failure retries. Un base model más pequeño con más inference search puede convenir para tareas de baja concurrencia y alto valor; un base model más potente puede reducir los pasos por solicitud. La elección debe basarse en la target workload y el SLO, no en qué etapa está de moda.
\end{knowledgebox}

\subsection{Resumen de la sección}

El frontier progress proviene de redistribuir el compute a lo largo del ciclo de vida, no de que una etapa sustituya de forma permanente a otra. El program debe elegir la combinación con base en la quality, el cost, la latency y el risk de extremo a extremo.
